% !TEX root = main.tex
\section{References}\label{references}


{[}1{]} Holt-Lunstad, J.; Smith, T.B.; Layton, J.B. Social Relationships and Mortality Risk: A Meta-Analytic Review. \emph{PLoS Med.} \textbf{2010}, \emph{7}, e1000316.

{[}2{]} Masterson Creber, R.; Rozenfeld, A.; Lucette, A.; et al.~Social Isolation in Older Adults: A Scoping Review. \emph{J. Am. Geriatr. Soc.} \textbf{2020}, \emph{68}, 1622--1634.

{[}3{]} Costa, P.T.; McCrae, R.R. \emph{Revised NEO Personality Inventory (NEO-PI-R) and NEO Five-Factor Inventory (NEO-FFI): Professional Manual}; Psychological Assessment Resources: Odessa, FL, USA, 1992.

{[}4{]} Roberts, B.W.; Kuncel, N.R.; Shiner, R.; Caspi, A.; Goldberg, L.R. The Power of Personality: The Comparative Validity of Personality Traits, Socioeconomic Status, and Cognitive Ability for Predicting Important Life Outcomes. \emph{Perspect. Psychol. Sci.} \textbf{2007}, \emph{2}, 313--345.

{[}5{]} Azaria, A.; Azoulay, R.; Reches, S. ChatGPT Is a Remarkable Tool for Experts. \emph{Nat. Med.} \textbf{2023}, \emph{29}, 1663--1664.

{[}6{]} Mairesse, F.; Walker, M.A.~Controlling User Perceptions of Linguistic Style: Trainable Generation of Personality Traits. \emph{Comput. Linguist.} \textbf{2011}, \emph{37}, 455--488.

{[}7{]} Wu, W.; Heierli, J.; Meisterhans, M.; Moser, A.; Färber, A.; Dolata, M.; Gavagnin, E.; de Spindler, A.; Schwabe, G. PROMISE: A Framework for Developing Complex Conversational Interactions. \emph{arXiv} \textbf{2024}, arXiv:2312.03699.

{[}8{]} Bickmore, T.W.; Picard, R.W. Establishing and Maintaining Long-Term Human-Computer Relationships. \emph{ACM Trans. Comput.-Hum. Interact.} \textbf{2005}, \emph{12}, 293--327.

{[}9{]} John, O.P.; Naumann, L.P.; Soto, C.J. Paradigm Shift to the Integrative Big Five Trait Taxonomy: History, Measurement, and Conceptual Issues. In \emph{Handbook of Personality: Theory and Research}, 3rd ed.; John, O.P., Robins, R.W., Pervin, L.A., Eds.; Guilford Press: New York, NY, USA, 2008; pp.~114--158.

{[}10{]} Quirin, M.; Malekzad, F.; Paudel, D.; Knoll, A.C.; Mirolli, M. Dynamics of Personality: The Zurich Model of Motivation Revived, Extended, and Applied to Personality. \emph{J. Pers.} \textbf{2023}, \emph{91}, 928--946.

{[}11{]} Krippendorff, K. \emph{Computing Krippendorff's Alpha-Reliability}; Departmental Papers (ASC), University of Pennsylvania: Philadelphia, PA, USA, 2011.

{[}12{]} Alisamir, S.; Ringeval, F. On the Evolution of Speech Representations for Affective Computing: A Brief History and Critical Overview. \emph{IEEE Signal Process. Mag.} \textbf{2021}, \emph{38}, 12--21.

{[}13{]} Alsharekh, M.F. Facial Emotion Recognition in Verbal Communication Based on Deep Learning. \emph{Electronics} \textbf{2022}, \emph{11}, 2568.

{[}14{]} Alberts, L.; Keeling, G.; McCroskery, A. Should Agentic Conversational AI Change How We Think about Ethics? Characterising an Interactional Ethics Centred on Respect. \emph{arXiv} \textbf{2024}, arXiv:2401.09187.

{[}15{]} Zheng, Z.; Liao, L.; Deng, Y.; Nie, L. Building Emotional Support Chatbots in the Era of LLMs. \emph{arXiv} \textbf{2023}, arXiv:2308.11584.

{[}16{]} Chen, K.; Kang, X.; Lai, X.; Ni, Z. Enhancing Emotional Support Capabilities of Large Language Models through Cascaded Neural Networks. In \emph{Proceedings of the 2023 4th International Conference on Computer, Big Data and Artificial Intelligence (ICCBD+AI)}; 2023; pp.~318--326.

{[}17{]} Dong, T.; Liu, F.; Wang, X.; Jiang, Y.; Zhang, X.; Sun, X. EmoAda: A Multimodal Emotion Interaction and Psychological Adaptation System. In \emph{Proceedings of the Conference on Multimedia Modeling}; Springer: Cham, Switzerland, 2024.

{[}18{]} Abbasian, M.; Azimi, I.; Rahmani, A.M.; Jain, R.C. Conversational Health Agents: A Personalized LLM-Powered Agent Framework. \emph{arXiv} \textbf{2023}, arXiv:2310.02293.

{[}19{]} Sorino, P.; et al.~ARIEL: Brain-Computer Interfaces Meet Large Language Models for Emotional Support Conversation. In \emph{Adjunct Proceedings of the 32nd ACM Conference on User Modeling, Adaptation and Personalization (UMAP '24)}, Cagliari, Italy, 1--4 July 2024.

{[}20{]} Dongre, P. Physiology-Driven Empathic Large Language Models (EmLLMs) for Mental Health Support. In \emph{Extended Abstracts of the CHI Conference on Human Factors in Computing Systems (CHI EA '24)}, Honolulu, HI, USA, 11--16 May 2024.

{[}21{]} Ta, V.; Griffith, C.; Boatfield, C.; Wang, X.; Civitello, M.; Maffei, H. User Experiences of Social Support from Companion Chatbots in Everyday Contexts: Thematic Analysis. \emph{J. Med. Internet Res.} \textbf{2020}, \emph{22}, e16235.

{[}22{]} Broadbent, E.; et al.~ElliQ, an AI-Driven Social Robot to Alleviate Loneliness: Progress and Lessons Learned. \emph{JAR Life} \textbf{2024}, \emph{13}, 22--28.

{[}23{]} Shah, S.G.S.; et al.~Effectiveness of Digital Technology Interventions to Reduce Loneliness in Adults: A Protocol for a Systematic Review and Meta-Analysis. \emph{BMJ Open} \textbf{2019}, \emph{9}, e029324.

{[}24{]} Shah, S.G.S.; et al.~Evaluation of the Effectiveness of Digital Technology Interventions to Reduce Loneliness in Older Adults: Systematic Review and Meta-Analysis. \emph{J. Med. Internet Res.} \textbf{2021}, \emph{23}, e24712.

{[}25{]} Zhang, H.; Chen, Y.; Wang, M.; Feng, S. FEEL: A Framework for Evaluating Emotional Support Capability with Large Language Models. \emph{arXiv} \textbf{2024}, arXiv:2403.15699.

{[}26{]} Berlyne, D.E. \emph{Conflict, Arousal, and Curiosity}; McGraw-Hill: New York, NY, USA, 1960.

{[}27{]} Hebb, D.O. \emph{The Organization of Behavior: A Neuropsychological Theory}; Wiley: New York, NY, USA, 1949.

{[}28{]} Bischof, N. \emph{Das Rätsel Ödipus: Die biologischen Wurzeln des Urkonfliktes}; Piper: Munich, Germany, 1985.

{[}29{]} Bischof, N. Untersuchungen zur Systemanalyse der sozialen Motivation I: Die Tantalus-Situation. \emph{Z. Psychol.} \textbf{1993}, \emph{201}, 5--43.

{[}30{]} Zhou, M.X.; et al.~Designing Effective Interview Chatbots: Automatic Chatbot Profiling and Design Suggestion Generation. In \emph{Proceedings of the CHI Conference on Human Factors in Computing Systems (CHI '21)}, Yokohama, Japan, 8--13 May 2021.

{[}31{]} Montgomery, D.C. \emph{Design and Analysis of Experiments}, 9th ed.; Wiley: Hoboken, NJ, USA, 2017.

{[}32{]} Collins, L.M.; Dziak, J.J.; Li, R. Design of Experiments with Multiple Independent Variables: A Resource Management Perspective on Complete and Fractional Factorial Designs. \emph{Psychol. Methods} \textbf{2009}, \emph{14}, 202--224.

{[}33{]} Carroll, C.; Patterson, M.; Wood, S.; Booth, A.; Rick, J.; Balain, S. A Conceptual Framework for Implementation Fidelity. \emph{Implement. Sci.} \textbf{2007}, \emph{2}, 40.

{[}34{]} Faul, F.; Erdfelder, E.; Buchner, A.; Lang, A.G. Statistical Power Analyses Using G*Power 3.1: Tests for Correlation and Regression Analyses. \emph{Behav. Res. Methods} \textbf{2009}, \emph{41}, 1149--1160.

{[}35{]} Schulz, K.F.; Altman, D.G.; Moher, D.; CONSORT Group. CONSORT 2010 Statement: Updated Guidelines for Reporting Parallel Group Randomised Trials. \emph{BMJ} \textbf{2010}, \emph{340}, c332.

{[}36{]} Chan, A.W.; Tetzlaff, J.M.; Gøtzsche, P.C.; Altman, D.G.; Mann, H.; Berlin, J.A.; Dickersin, K.; Hróbjartsson, A.; Schulz, K.F.; Parulekar, W.R.; et al.~SPIRIT 2013 Statement: Defining Standard Protocol Items for Clinical Trials. \emph{Ann. Intern. Med.} \textbf{2013}, \emph{158}, 200--207.

{[}37{]} World Medical Association. World Medical Association Declaration of Helsinki: Ethical Principles for Medical Research Involving Human Subjects. \emph{JAMA} \textbf{2013}, \emph{310}, 2191--2194.

{[}38{]} D'Alfonso, S. AI in Mental Health. \emph{Curr. Opin. Psychol.} \textbf{2020}, \emph{36}, 112--117.

{[}39{]} Torous, J.; Myrick, K.J.; Rauseo-Ricupero, N.; Firth, J. Digital Mental Health and COVID-19: Using Technology Today to Accelerate the Curve on Access and Quality Tomorrow. \emph{JMIR Ment. Health} \textbf{2020}, \emph{7}, e18848.

{[}40{]} McCrae, R.R.; Costa, P.T., Jr.~\emph{Personality in Adulthood: A Five-Factor Theory Perspective}, 2nd ed.; Guilford Press: New York, NY, USA, 2003.

{[}41{]} McCrae, R.R.; John, O.P. An Introduction to the Five-Factor Model and Its Applications. \emph{J. Pers.} \textbf{1992}, \emph{60}, 175--215.

{[}42{]} Goldberg, L.R. The Structure of Phenotypic Personality Traits. \emph{Am. Psychol.} \textbf{1993}, \emph{48}, 26--34.

{[}43{]} Gelman, A.; Carlin, J.B.; Stern, H.S.; Dunson, D.B.; Vehtari, A.; Rubin, D.B. \emph{Bayesian Data Analysis}, 3rd ed.; CRC Press: Boca Raton, FL, USA, 2013.

{[}44{]} Funder, D.C. On the Accuracy of Personality Judgment: A Realistic Approach. \emph{Psychol. Rev.} \textbf{1995}, \emph{102}, 652--670.

{[}45{]} Vazire, S. Who Knows What about a Person? The Self--Other Knowledge Asymmetry (SOKA) Model. \emph{J. Pers. Soc. Psychol.} \textbf{2010}, \emph{99}, 281--303.

{[}46{]} OpenAI. \emph{GPT-4 Technical Report}; arXiv:2303.08774, 2023.

{[}47{]} Brown, T.; Mann, B.; Ryder, N.; Subbiah, M.; Kaplan, J.D.; Dhariwal, P.; Neelakantan, A.; Shyam, P.; Sastry, G.; Askell, A.; et al.~Language Models are Few-Shot Learners. In \emph{Proceedings of the 34th International Conference on Neural Information Processing Systems (NIPS '20)}, Vancouver, BC, Canada, 6--12 December 2020; pp.~1877--1901.

{[}48{]} Reynolds, L.; McDonell, K. Prompt Programming for Large Language Models: Beyond the Few-Shot Paradigm. In \emph{Extended Abstracts of the 2021 CHI Conference on Human Factors in Computing Systems}, Yokohama, Japan, 8--13 May 2021; pp.~1--7.

{[}49{]} White, J.; Fu, Q.; Hays, S.; Sandborn, M.; Olea, C.; Gilbert, H.; Elnashar, A.; Spencer-Smith, J.; Schmidt, D.C. A Prompt Pattern Catalog to Enhance Prompt Engineering with ChatGPT. \emph{arXiv} \textbf{2023}, arXiv:2302.11382.

{[}50{]} Liu, P.; Yuan, W.; Fu, J.; Jiang, Z.; Hayashi, H.; Neubig, G. Pre-train, Prompt, and Predict: A Systematic Survey of Prompting Methods in Natural Language Processing. \emph{ACM Comput. Surv.} \textbf{2023}, \emph{55}, 1--35.

{[}51{]} APA. \emph{Practice Guideline for the Treatment of Patients with Major Depressive Disorder}, 3rd ed.; American Psychiatric Association: Arlington, VA, USA, 2010.

{[}52{]} Simon, G.E. Patient-Centered Outcomes in Mental Health Care. \emph{JAMA} \textbf{2011}, \emph{306}, 1141--1142.

{[}53{]} McCrae, R.R.; Costa, P.T., Jr.~\emph{Revised NEO Personality Inventory (NEO PI-R) and NEO Five-Factor Inventory (NEO-FFI) Professional Manual}; Psychological Assessment Resources: Odessa, FL, USA, 1992.

{[}54{]} Likert, R. A Technique for the Measurement of Attitudes. \emph{Arch. Psychol.} \textbf{1932}, \emph{22}, 5--55.

{[}55{]} Zheng, L.; Chiang, W.L.; Sheng, Y.; Zhuang, S.; Wu, Z.; Zhuang, Y.; Lin, Z.; Li, Z.; Li, D.; Xing, E.; et al.~Judging LLM-as-a-judge with MT-Bench and Chatbot Arena. \emph{arXiv} \textbf{2023}, arXiv:2306.05685.

{[}56{]} Kim, S.; Lee, S.; Kim, J.; Yoo, H.; Seol, Y.Y.; Kim, S.; Cho, H.; Choi, S.B.; Kim, S.B. Evaluating LLM-as-a-judge with Professional Human Ratings. \emph{arXiv} \textbf{2024}, arXiv:2402.18139.

{[}57{]} Landis, J.R.; Koch, G.G. The Measurement of Observer Agreement for Categorical Data. \emph{Biometrics} \textbf{1977}, \emph{33}, 159--174.

{[}58{]} Gisev, N.; Bell, J.S.; Chen, T.F. Interrater Reliability and Agreement in Clinical Research: A Guide to Best Practice. \emph{Res. Social Adm. Pharm.} \textbf{2013}, \emph{9}, 330--337.

{[}59{]} Little, R.J.; Rubin, D.B. \emph{Statistical Analysis with Missing Data}, 3rd ed.; Wiley: Hoboken, NJ, USA, 2019.

{[}60{]} Cohen, J. \emph{Statistical Power Analysis for the Behavioral Sciences}, 2nd ed.; Lawrence Erlbaum Associates: Hillsdale, NJ, USA, 1988.

{[}61{]} Efron, B.; Tibshirani, R.J. \emph{An Introduction to the Bootstrap}; Chapman and Hall/CRC: New York, NY, USA, 1994.

{[}62{]} Wasserstein, R.L.; Lazar, N.A. The ASA's Statement on p-Values: Context, Process, and Purpose. \emph{Am. Stat.} \textbf{2016}, \emph{70}, 129--133.

{[}63{]} Chen, L.; Zaharia, M.; Zou, J. How is ChatGPT's Behavior Changing over Time? \emph{arXiv} \textbf{2023}, arXiv:2307.09009.

{[}64{]} Ouyang, L.; Wu, J.; Jiang, X.; Almeida, D.; Wainwright, C.; Mishkin, P.; Zhang, C.; Agarwal, S.; Slama, K.; Ray, A.; et al.~Training Language Models to Follow Instructions with Human Feedback. In \emph{Proceedings of the 36th International Conference on Neural Information Processing Systems (NIPS '22)}, New Orleans, LA, USA, 28 November--9 December 2022; pp.~27730--27744.

{[}65{]} Ahmad, R.; Siemon, D.; Gnewuch, U.; Robra-Bissantz, S. Designing Personality-Adaptive Conversational Agents for Mental Health Care. \emph{Inf. Syst. Front.} \textbf{2022}, \emph{24}, 923--943.

{[}66{]} Wanniarachchi, V.U.; Greenhalgh, C.; Choi, A.; Warren, J.R. Personalization Variables in Digital Mental Health Interventions for Depression and Anxiety in Adolescents and Youth: A Scoping Review. \emph{Front. Digit. Health} \textbf{2025}, \emph{7}, 1500220.

{[}67{]} Soni, U.; Thakkar, H.; Pandya, A.; Patel, D. The Power of Personalization: A Systematic Review of Personality-Adaptive Chatbots. \emph{SN Comput. Sci.} \textbf{2023}, \emph{4}, 661.

{[}68{]} Devdas, S. \emph{Enhancing Emotional Support through Conversational AI via Big Five Personality Detection and Behavior Regulation Based on the Zurich Model}; Master's Thesis, Lucerne University of Applied Sciences and Arts: Luzern, Switzerland, 2025.

